\documentclass[parskip=half-]{scrartcl}
\usepackage[french]{babel}
\usepackage[T1]{fontenc}
\usepackage[utf8]{inputenc}
\usepackage{mathtools}
\usepackage{tabularray}
\usepackage{enumitem}
\usepackage{graphicx}
\usepackage[charter]{mathdesign}
\begin{document}
% \section*{Matériel}
% Chaque enveloppe (le jeu est conçu pour être par 5 équipes) doit contenir les indices qui permettront de résoudre les différentes énigmes.
% \begin{enumerate}
%     \item 5 jeux de picross. Chaque jeu contient 5 picross.
%     \item 5 feuilles de règle du picross.
%     \item 5 photos représentant le kidnapping du nombre pi. 
%     \item 5 jeux de photos de pizzeria. Chaque jeu contient 6 pizzerias.
%     \item 5 QRcodes à compléter.
%     \item 5 jeux d'étiquettes. Chaque jeu contient 3 étiquettes.
% \end{enumerate}

\section*{Avant de jouer}
Les participants se répartissent équitablement en 5 équipes. Ils se munissent d'un bic ou d'un crayon et d'une calculatrice.

\section*{Énigme 1}
Chaque équipe reçoit un jeu de picross contenant 5 picross ainsi qu'une feuille de règle. Ils doivent colorier sans se tromper des cases en noir afin de que les cases restantes (en blanc) forment un chiffre (1,3,4 ou 5). Une fois ces 5 chiffres obtenus, les participants doivent les réarranger pour former la combinaison suivante : 31415. La première équipe qui termine l'énigme y mets fin.

\section*{Énigme 2}
Chaque équipe reçoit une photo représentant le kidnapping du nombre pi, une bande de papier transparent de couleur rouge et un jeu de photos représentant des pizzerias.
À l'aide du filtre rouge, un indice est décodé (5 indices au total dont un sans filtre rouge). Les équipes doivent échanger leurs indices (en les écrivant par exemple sur un tableau) qui sont en fait les caractéristiques de la pizzeria dans laquelle pi a été enlevé. En procédant par élimination, les équipes doivent arriver à un consensus et ne garder qu'une seule pizzeria.

\section*{Énigme 3}
Chaque équipe reçoit une note sur laquelle figure un message ainsi qu'un QRCode (les QRCodes doivent être distribués par rapport à la bonne photo, cf. numéro de la photo en bas à droite.). Les équipes doivent colorier sur le QRCode les deux nombres dissimulés dans la photo du kidnapping. En utilisant un smartphone (le leur ou celui d'un des meneurs de jeu), chaque équipe trouvera un bout de phrase. Une fois la phrase complétée ("Le code est donné par les chiffres qui se situent entre la troisième et la huitième décimale de pi") il découvre le code pour entrer dans la pizzeria : 5926. Ce code pourra être entré sur le digicode papier en y plaçant correctement des aimants.

\section*{Énigme 4}
Pendant le déroulement des autres énigmes, le meneur de jeu aura dissimulé discrètement la clé d'un cadenas dans la classe. Une fois la clé retrouvée elle permettra de déverrouiller le cadenas qui verrouille une armoire ou une boîte. À l'intérieur de celle-ci, se trouve une autre boîte verrouillée par un cadenas à code.

\section*{Énigme 5}
Chaque équipe reçoit 3 étiquettes qui représentent un cercle et son rayon. La longueur du rayon est donnée. Le but est de calculer le périmètre du cercle, l'aire du disque ou le volume de la sphère et de ne conserver que les chiffres des dizaines et des unités. Ceux-ci sont alors inscrits dans un tableau (tableur de préférence). La somme de nombres ainsi obtenus donne le code du cadenas.
\end{document}